% ============================================================
%  ch02a  §2 引言 + §2.1 范畴算术局部朗兰兹
%  原文对应：ICM-arXiv.tex 782--935 行
%  编号：定理 2.1.1；定理 2.1.2；猜想 2.1.3；注 2.1.4；
%        公式 (2.1)(2.2)(2.3)(2.4)；定理 2.1.5；例 2.1.6；例 2.1.7；注 2.1.8
% ============================================================

\section{从几何到经典朗兰兹纲领}

上一节讨论了范畴化与几何化的想法如何引出几何朗兰兹纲领的发展。另一方面，量子物理的想法允许我们在不同维数的流形上取值（拓扑）量子场论，从而把 \eqref{categorification} 中的箭头反向。这类想法在几何与拓扑中当然并不新鲜；但出人意料的是，它也带来了对经典朗兰兹纲领的新理解。事实上，众所周知整体域与 $3$-流形之间存在类比，在此类比下 Frobenius 对应于圆周的基本群。于是“场论在圆周上的紧化”引出范畴迹方法（例如见 \cite{AGKRRV, BCHN, Zhu-CDM}），它在几何表示论中起着越来越重要的作用。

\subsection{范畴算术局部朗兰兹}\label{S: Local Langlands}

本小节中，设 $F$ 是 $\bQ_p$ 的有限扩张，或同构于 $\bF_q((\varpi))$。经典局部朗兰兹对应寻求用 Galois 数据来分类 $G(F)$ 的光滑不可约表示。除 $G=\GL_n$ 情形（这是 \cite{LRS-GLn, Harris-Taylor-GLn} 的定理）之外，其精确表述相当复杂。然而，局部几何朗兰兹是 2-范畴式的这一信条（见注 \ref{2-categorical LGL}）提示：即便经典局部朗兰兹对应也应当、而且多半需要被范畴化。

表述范畴算术局部朗兰兹所需的第一个要素是下述结果，它由 \cite{DHKM, Fargues-Scholze-Geometrization, Zhu-Conjecture} 各自独立得到。我们采用 \cite{Zhu-Conjecture} 中的表述，其中（强）连续同态的概念参见该文。

\begin{thm}\label{Stack of local parameters}
把 $\bZ_\ell$-代数 $A$ 送到满足 $d\circ\rho=(\cycl^{-1},\pr)$ 的（强）连续同态 $\rho: W_F\to \CG(A)$ 所构成空间的预叠，可由一个（经典）概形 $\Loc_{\CG}^\Box$ 表示；该概形是若干仿射概形的无交并，每个仿射概形在 $\bZ_\ell$ 上平坦、有限型，且局部完全交。
\end{thm}

$\hat{G}$ 在 $\CG$ 上的共轭作用诱导出 $\hat{G}$ 在 $\Loc_{\CG}^\Box$ 上的作用；称商叠 $\Loc_{\CG}=\Loc_{\CG}^\Box/\hat{G}$ 为局部朗兰兹参数叠，它粗略地说把上述诸 $\rho$ 构成的群胚按 $\hat{G}$-共轭分类。

在局部朗兰兹对应的范畴版本中，Galois 一侧自然要考虑 $\Loc_{\CG}$ 上凝聚层构成的（$\infty$-）范畴 $\Coh(\Loc_{\CG})$。表示一侧，人们或许会朴素地考虑 $G(F)$ 光滑表示构成的（$\infty$-）范畴 $\Rep(G(F),\La)$。但实际上这个范畴需要扩大。

这可以从不同角度看出。事实上，许多人共有的一条普遍经验是：在经典局部朗兰兹对应中，最好把 $G$ 的表示与一族群（它们是 $G$ 的内形式的精细化）放在一起研究。这类群的族有各种提法。算术几何（即对志村簇与 Shtuka 模空间的 $p$-进与 mod $p$ 几何的研究）与几何表示论（即范畴迹构造）都提示，应当研究由 Kottwitz 集
\[
B(G) = G(\breve F)/\sim, \quad \mbox{其中 } g\sim g' \mbox{ 指存在 } h\in G(\breve F) \mbox{ 使 } g'=h^{-1}g\sigma(h),
\]
所给出的一族群 $\{J_b(F)\}_{b\in B(G)}$ 的表示范畴粘合而成的范畴。这里对 $b\in B(G)$（提升为 $G(\breve F)$ 中的元素），群 $J_b$ 是按下述方式定义的 $F$-群：对 $F$-代数 $R$，令 $J_b(R)=\{h\in G(\breve F\otimes_FR)\mid h^{-1}b\sigma(h)=b\}$。特别地，若 $b=1$ 则 $J_b=G$。一般地，存在良定义到共轭的嵌入 $(J_{b})_{\overline{F}}\to G_{\overline{F}}$，它使 $J_b$ 成为 $G$ 的某个 Levi 子群的内形式的精细化（设 $G$ 拟分裂）。

有两种方式把这想法精确化。一种属于 Fargues--Scholze \cite{Fargues-Scholze-Geometrization}：他们把 $B(G)$ 看作 Fargues--Fontaine 曲线上 $G$-丛的 $v$-叠 $\Bun_G$ 的点集，并考虑 $\Bun_G$ 上适当定义的 étale 层构成的范畴 $D_{\on{lis}}(\Bun_G,\La)$，后者确实把所有 $\Rep(J_b(F),\La)$ 粘合在一起。需要指出，这一途径依赖 Scholze 关于 diamond 与凝聚态数学的 $\ell$-进形式体系的工作。

另一途径 \cite{Xiao-Zhu-Cycle, Zhu-CDM, Zhu-Conjecture, Hemo-Zhu-Unipotent} 与范畴迹的想法密切相关：它把 $B(G)$ 看作（étale）商叠
\begin{equation*}\label{E: kot stack}
\frakB(G):=LG/\Ad_\sigma LG
\end{equation*}
的点集，其中 $\Ad_\sigma$ 表示由
\[
\Ad_\sigma: LG\times LG\to LG,\quad  (h,g)\mapsto hg\sigma(h)^{-1}
\]
给出的 Frobenius 扭共轭。于是有前面提到的 $\La$-层范畴 $\Shv(\kot(G)\otimes \bar{k},\La)$。尽管 $\kot(G)$ 在传统代数几何中是病态对象，关于它的几何仍可以说上几句，而范畴 $\Shv(\frakB(G)\otimes \bar{k},\La)$ 相当合理。此外，还可以定义 $\frakB(G)\otimes \bar{k}$ 上可构造层构成的范畴 $\Shv_c(\frakB(G)\otimes \bar{k},\La)$，下面略作说明，细致讨论见 \cite{Hemo-Zhu-Unipotent}。

对每个 $k$ 上的代数闭域 $\Omega$，$\kot(G)$ 的 $\Omega$-点群胚就是 $\Omega$ 上带 $G$-结构的 $F$-同晶体的群胚，其同构类集合同于 Kottwitz 集 $B(G)$。然而 $\kot(G)$ 并不只是其各点的无交并；它带有一个分层，称为 Newton 分层，以 $B(G)$ 标号。具体地说，$B(G)$ 上有自然的偏序，粗略地讲，对每个 $b\in B(G)$，对应于 $b'\leq b$ 的那些 $\Omega$-点构成闭子叠 $i_{\leq b}: \kot(G)_{\leq b}\subset \kot(G)\otimes \bar k$，而对应于 $b$ 的那些 $\Omega$-点构成开子叠 $j_{b}: \kot(G)_b\subset \kot(G)_{\leq b}$。

特别地，$B(G)$ 中的基本元（即关于偏序 $\leq$ 的极小元）给出闭层。还要指出，若 $b$ 是基本元，则 $J_b$ 是 $G$ 的某个内形式的精细化，通常称为 $G$ 的扩张纯内形式。

本小节余下部分把 $\kot(G)\otimes \bar k$ 简记为 $\kot(G)$，并记 $i_b=i_{\leq b}j_b: \kot(G)_b\hookrightarrow \kot(G)$ 为局部闭嵌入。对 $b$，以 $\fgRep(J_b(F),\La)$ 记 $\Rep(J_b(F),\La)$ 中由开紧子群 $K\subset J_b(F)$ 的平凡表示紧诱导出的表示
\[
\delta_{K,\La}:=c\on{-ind}_K^{J_b(F)}(\La)
\]
在有限余极限与收缩下生成的完全子范畴。下述定理取自 \cite{Hemo-Zhu-Unipotent}，它概括了 $\Shv_c(\kot(G),\La)$ 的若干性质。

\begin{thm}\label{T: Local-Langlands-Category}
\begin{enumerate}
    \item $\Shv(\kot(G),\La)$ 中的对象可构造，当且仅当它对每个 $\frakB(G)_b$ 的 $!$-限制可构造，且对几乎所有 $b$ 为零。若 $\La$ 是特征零的域，则 $\Shv_{c}(\frakB(G),\La)$ 由 $\Shv(\kot(G),\La)$ 的紧对象组成。
    \item 对每个 $b\in B(G)$，存在典范等价 $\Shv_c(\frakB(G)_b,\La)\cong \fgRep(J_b(F),\La)$。存在完全忠实嵌入 $i_{b,*}, i_{b,!}:  \Shv_c(\frakB(G)_b,\La)\to \Shv_c(\kot(G),\La)$（当 $b$ 为基本元时两者相同），它们给出 $\Shv_c(\frakB(G),\La)$ 关于 $\{\Shv_c(\frakB(G)_b,\La)\}_b$ 的半正交分解。
    \item 存在自对偶函子 $\bD^{\coh}\colon \Shv_c(\frakB(G),\La)\simeq \Shv_c(\frakB(G),\La)^{\vee}$，其构造方式是把各个 $\fgRep(J_b(F),\La)$ 上的上同调对偶（Bernstein--Zelevinsky 意义下）粘合起来。
    \item 把各个 $\fgRep(J_b(F),\La)$ 上（平移后的）$t$-结构粘合起来，得到自然的反常 $t$-结构；当 $\La$ 是域时，它被 $\bD^{\coh}$ 保持。
\end{enumerate}
\end{thm}

下述算术局部朗兰兹猜想的范畴形式（\cite[Sect. 4.6]{Zhu-Conjecture}）受到 \S\ref{Global GLC} 中讨论的整体几何朗兰兹猜想的启发。

\begin{conj}\label{C: LLC}
设 $G$ 在 $F$ 上拟分裂，并带有钉 $(B,T,e)$；固定非平凡加性特征 $\psi: F\to \bZ_\ell[\mu_{p^\infty}]^\times$。则存在范畴之间的典范等价
\[
\bL_{G}:   \Coh(\Loc_{\CG}\otimes\La)\cong \Shv_{c}(\frakB(G),\La).
\]
\end{conj}

\begin{rmk}
(1) 该猜想有一个密切相关的版本：把 $\Shv_{c}(\frakB(G),\La)$ 换成 $\Shv(\frakB(G),\La)$，把 $\Coh(\Loc_{\CG}\otimes\La)$ 换成其 ind-完备化（并加上某些支撑条件）（见 \cite[Sect. 4.6]{Zhu-Conjecture}）。Fargues--Scholze \cite{Fargues-Scholze-Geometrization} 提出了一个与此版本平行的猜想，其中范畴 $\Shv(\frakB(G),\La)$ 被换成了前面提到的 $D_{\on{lis}}(\Bun_G,\La)$。

(2) \cite{Zhu-Conjecture} 还说明了在 $\bQ$ 上建立此类等价的一个 motivic 期望。
\end{rmk}

该猜想的一个推论是：对每个 $b$，应当存在完全忠实嵌入
\begin{equation*}\label{E: functor A}
\frakA_{J_b}:  \fgRep(J_b(F),\La)\to  \Coh(\Loc_{\CG}\otimes\La),
\end{equation*}
它是把 $\bL_G$ 的拟逆限制到 $i_{b,!}(\fgRep(J_b(F),\La))$ 而得到的。这样的函子之存在，与族中局部朗兰兹的想法密切相关；在 $J_b=G$ 分裂且 $\La$ 是特征零的域的情形，它也在 \cite{Hellmann-Iwahori, BCHN} 中被考虑过。

特别地，对每个开紧子群 $K\subset J_b(F)$，应当存在 $\Loc_{\CG}\otimes\La$ 上的凝聚层
\begin{equation}\label{E: sheaf A}
\frakA_{K,\La}:=\frakA_{J_b}(\delta_{K,\La}),
\end{equation}
使得
\begin{equation}\label{E: Hecke alg isom}
\big(R\End_{\Coh(\Loc_{\CG}\otimes\La)}\frakA_{K,\La}\big)^{\op}\cong \big(R\End_{\Rep(G(F),\La)}\delta_{K,\La}\big)^{\op}=:H_{K,\La}.
\end{equation}
代数 $H_{K,\La}$ 有时称为导出 Hecke 代数，因为（当 $\La=\bZ_\ell$ 或 $\bF_\ell$ 时）它未必集中在上同调次数零。$\frakA_{K,\La}$ 在各种情形下的猜想性描述见 \cite[Sect. 4.3-4.5]{Zhu-Conjecture}。

与整体几何朗兰兹猜想一样，猜想 \ref{C: LLC} 中的等价应当满足一族相容性条件。例如，它应当与两侧的抛物诱导相容，也应当与 $\Shv_c(\frakB(G),\La)$ 上的对偶 $\bD^{\coh}$ 以及 $\Coh(\Loc_{\CG}\otimes\La)$ 的（修正的）Grothendieck--Serre 对偶相容。细节见 \cite{Zhu-Conjecture, Hemo-Zhu-Unipotent}。

另一方面，猜想 \ref{C: LLC} 预测完全复形的范畴 $\Perf(\Loc_{\CG}\otimes\La)$ 作用于 $\Shv_c(\frakB(G),\La)$，这类似于 \S\ref{Global GLC} 中提到的谱作用。\cite{Fargues-Scholze-Geometrization} 的主要结果之一，就是在他们的框架中构造出这样的作用。目前人们还不知道 $\Shv_c(\frakB(G),\La)$ 上是否存在这样的谱作用，但有一些令人信服的证据表明猜想 \ref{C: LLC} 仍应成立。

如前面一样，设 $G$ 可延拓为 $\mO$ 上的约化群。于是有闭子叠
\[\Loc_{\CG}^{\ur}\subset\Loc_{\CG}^{\un}\subset \Loc_{\CG},\]
通常分别称为非分歧参数叠与幂幺参数叠，它们分类那些满足 $\rho(I_F)$ 平凡（分别地，$\rho(I_F)$ 幂幺）的 $\rho$。当 $\La=\bQ_\ell$ 时，$\Loc_{\CG}^{\un}\otimes\bQ_\ell$ 是 $\on{Loc}_{\CG}\otimes\bQ_\ell$ 的一个连通分支。

另一方面，有幂幺子范畴 $\Shv_c^{\un}(\kot(G),\bQ_\ell)\subset \Shv_c(\kot(G),\bQ_\ell)$，粗略地说它是所有 $b\in B(G)$ 的 $J_b(F)$ 幂幺表示范畴 $\fgRep^{\un}(J_b(F),\bQ_\ell)$（由 \cite{Lusztig-Unipotent} 引入）的粘合。下述定理取自 \cite{Hemo-Zhu-Unipotent}，由对定理 \ref{T: Bezrukavnikov equivalence} 取 Frobenius 扭的范畴迹得到。

\begin{thm}\label{T: tame local Langlands}
设 $G$ 是 $\mO$ 上的约化群，并固定钉 $(B,T,e)$。则存在典范等价
\[
\bL^{\un}_G:  \Coh(\Loc_{\CG}^{\un}\otimes\bQ_\ell)\cong  \Shv_{c}^{\un}(\frakB(G),\bQ_\ell).
\]
\end{thm}

对算术应用而言，在此等价下把具体对象对应起来是重要的。我们给出几个例子，更多的对应见 \cite{Hemo-Zhu-Unipotent}（并见 \cite[Sect. 4.3-4.5]{Zhu-Conjecture}）。

\begin{ex}
等价 $\bL_G^{\un}$ 对所有抛孔子群 $K\subset G(F)$（Bruhat--Tits 意义下）给出 \eqref{E: sheaf A} 中的猜想凝聚层，使得 \eqref{E: Hecke alg isom} 成立。例如 $\frakA_{G(\mO),\bQ_\ell}\cong \mO_{\Loc_{\CG}^{\ur}}\otimes\bQ_\ell$，由此得到
\begin{equation}\label{E: derived Satake}
\big(R\End_{\Coh(\Loc_{\CG})}\mO_{\Loc_{\CG}^{\ur}}\big)^{\op}\otimes\bQ_\ell\cong \big(R\End\ \delta_{G(\mO),\bQ_\ell}\big)^{\op}= H_{G(\mO),\bQ_\ell}\cong H_{G(\mO)}^{\cl}\otimes\bQ_\ell.
\end{equation}
由于 $\Loc_{\CG}^{\ur}\cong (\CG|_{d=(q,\sigma)})/\hat{G}$，取零次上同调就还原出 Satake 同构 \eqref{Sat isom}。此外，它还蕴含左端没有高阶上同调，而这一点并不显然。需要指出，\cite[Sect. 4.3]{Zhu-Conjecture} 猜测 $\frakA_{G(\mO),\bZ_\ell}\cong \mO_{\Loc_{\CG}^{\ur}}$，因此 \eqref{E: derived Satake} 中第一个同构应当在 $\bZ_\ell$ 上成立，这称为（猜想性的）导出 Satake 同构。（但一般地 $H_{G(\mO),\bZ_\ell}\neq H_{G(\mO)}^{\cl}\otimes\bZ_\ell$。）

对 $\frakA_{I,\bQ_\ell}$ 还有一个纯 Galois 一侧的描述，即 \cite{BCHN, Zhu-Conjecture} 中定义的幂幺凝聚 Springer 层（并见 \cite{Hellmann-Iwahori}）。
\end{ex}

\begin{ex}\label{Ex: trace of Satake}
由构造，存在 $\bZ_\ell$ 上的叠的态射 $\Loc_{\CG}\to \bB\hat{G}$。设 $\hat{G}$ 在有限投射 $\La$-模上有表示，把它看作 $\bB\hat{G}_\La$ 上的向量丛，以 $\widetilde V$ 记其拉回至 $\Loc_{\CG}\otimes\La$；对 $?=\ur$ 或 $\un$，以 $\widetilde{V}^{?}\in \Perf(\Loc_{\CG}^{?}\otimes\La)$ 记它到 $\Loc_{\CG}^?\otimes\La$ 的限制。注意当 $\La=\bQ_\ell$ 时 $\widetilde{V}^{\ur}\cong \widetilde{V}\otimes \frakA_{G(\mO),\bQ_\ell}$。我们有
\[
\bL_G^{\un}(\widetilde{V}^{\ur})\cong \on{Nt}_!r^!\Sat(V)=:\mS_V,
\]
其中 $r$ 与 $\on{Nt}$ 是下述对应中的态射：
\[
\Hk_G=L^+G\backslash LG/L^+G\xleftarrow{r} LG/\Ad_\sigma L^+G\xrightarrow{\on{Nt}} LG/\Ad_\sigma LG=\kot(G).
\]
特别地，对 $\hat{G}$ 的两个表示 $V$ 与 $W$，存在与复合相容的态射
\begin{equation}\label{E: local S-operator}
R\Hom_{\Loc_{\CG}^{\ur}\otimes\bQ_\ell}(\widetilde{V}^{\ur}, \widetilde{W}^{\ur})\to R\Hom_{\Shv_c(\kot(G),\bQ_\ell)}(\mS_V,\mS_W).
\end{equation}
这样的映射最初由 \cite{Xiao-Zhu-Cycle,Zhu-CDM} 构造，其（非导出 $\Hom$ 空间的）版本随后在 \cite{Yu-Jacquet-Langlands-Higher-Weights} 中被推广到 $\bZ_\ell$ 系数。它有重要的算术应用，将在 \S\ref{Application} 中说明。
\end{ex}

\begin{rmk}
定理 \ref{T: tame local Langlands} 多半可以通过对定理 \ref{T: Monodromic Bezrukavnikov equivalence} 中的等价取 Frobenius 扭的范畴迹，推广到驯顺层。另一方面，如注 \ref{R: int Bez} 所述，把这些等价推广到 $\bZ_\ell$ 系数是重要的。
\end{rmk}
